% Work-MAT | 2. rocnik | hodina 16
\documentclass[11pt,a4paper]{article}
\input{../../sablony/pripravy/priprava.sty}
\usepackage{tikz}
\usetikzlibrary{arrows.meta}
\nastavPripravu{2. ročník}{Exponenciální rovnice}{16}{Exponenciální rovnice I: stejné základy}
\begin{document}
\poznamka{Navazujeme na prostotu exponenciální funkce z hodin 11--14, nikoli na logaritmy. Orientačně 7 min grafická motivace, 12 min hlavní pravidlo, 16 min příklady a 10 min samostatné upevnění.}
\blok{Co znamená neznámá v exponentu?}
Rovnice $2^x=8$ hledá průsečík grafu $y=2^x$ s přímkou $y=8$. Zápis $8=2^3$ ukazuje řešení $x=3$. Protože $2^x$ přísně roste, žádné jiné řešení mít nemůže.
\begin{definiceBox}[title={Definice - Exponenciální rovnice}]
\textbf{Exponenciální rovnicí} nazýváme rovnici, v níž se neznámá vyskytuje v exponentu mocniny. V této hodině řešíme hlavně rovnice, které lze převést na mocniny se stejným kladným základem různým od jedné.
\end{definiceBox}
\begin{zadaniBox}[title={Věta - Rovnost mocnin stejného základu}]
Pro $a>0$ a $a\ne1$ platí $a^{u(x)}=a^{v(x)}\iff u(x)=v(x)$, pokud jsou oba exponenty definované. Důvodem je prostota funkce $a^x$.
\end{zadaniBox}
\begin{prikladReseniBox}[title={Příklad 1}]
Vyřeš $3^{2x-1}=3^{x+2}$.
\tcblower
\[\begin{aligned}3^{2x-1}&=3^{x+2}\\2x-1&=x+2\\x&=3.\end{aligned}\]
Zkouška: $3^{2\cdot3-1}=3^5=3^{3+2}$. Množina kořenů $K=\{3\}$.
\end{prikladReseniBox}
\begin{prikladReseniBox}[title={Příklad 2}]
Vyřeš $2^{x+1}=\frac{1}{16}$.
\tcblower
\[\begin{aligned}2^{x+1}&=2^{-4}\\x+1&=-4\\x&=-5.\end{aligned}\]
Zkouška: $2^{-4}=\frac{1}{16}$. $K=\{-5\}$.
\end{prikladReseniBox}
\begin{prikladReseniBox}[title={Příklad 3}]
Rozhodni, zda má rovnice $5^{x-2}=-3$ reálné řešení.
\tcblower
Nemá, protože kladná mocnina $5^{x-2}$ nemůže být záporná. Tedy $K=\varnothing$. \emph{Záporný exponent není záporná funkční hodnota.}
\end{prikladReseniBox}
\poznamka{Chybu $a^u=a^v\Rightarrow u=v$ pro základ $a=1$ oprav protipříkladem $1^2=1^9$. Důležité jsou dvě podmínky: společný základ a $a>0$, $a\ne1$.}
\begin{zadaniReseniBox}[title={Samostatná práce 1}]
\begin{enumerate}[label=\alph*)]
\item $4^{x-1}=4^3$.
\item $\left(\frac{1}{2}\right)^{2x+1}=\left(\frac{1}{2}\right)^5$.
\item $7^{x+2}=1$.
\item $3^x=0$.
\item $5^{3x-2}=5^{x+6}$.
\item $\left(\frac{1}{3}\right)^{2x-3}=\left(\frac{1}{3}\right)^{x+2}$.
\item $9^{x^2-3x}=9^0$.
\item $2^{x^2-1}=2^{2x+2}$.
\end{enumerate}
\tcblower
\begin{enumerate}[label=\alph*)]
\item $x=4$.
\item $2x+1=5$, tedy $x=2$.
\item $x+2=0$, tedy $x=-2$.
\item $K=\varnothing$; $3^x>0$.
\item $3x-2=x+6$, tedy $x=4$.
\item $2x-3=x+2$, tedy $x=5$.
\item $x^2-3x=0$, tedy $K=\{0,3\}$.
\item $x^2-1=2x+2$, tedy $(x-3)(x+1)=0$ a $K=\{-1,3\}$.
\end{enumerate}
\end{zadaniReseniBox}
\begin{zadaniReseniBox}[title={Rozšíření 1}]
\begin{enumerate}[label=\alph*)]
\item Vyřeš $\left(\frac{1}{3}\right)^{x^2-4}=\left(\frac{1}{3}\right)^{x+2}$.
\item Vysvětli, proč i klesající exponenciální funkce dovoluje porovnání exponentů.
\end{enumerate}
\tcblower
\begin{enumerate}[label=\alph*)]
\item $x^2-4=x+2$ dává $(x-3)(x+2)=0$, tedy $K=\{-2,3\}$.
\item Je prostá: stejné hodnoty mají jen stejné argumenty.
\end{enumerate}
\end{zadaniReseniBox}
\begin{zadaniReseniBox}[title={Rozšíření 2}]
\begin{enumerate}[label=\alph*)]
\item Vyřeš rovnici $4^{\frac{x+1}{x-2}}=4^{\frac{x-3}{x+1}}$ s lineárně lomenými exponenty.
\item Uveď podmínky a zkontroluj, zda je výsledek přípustný.
\end{enumerate}
\tcblower
\begin{enumerate}[label=\alph*)]
\item Podmínky: $x\ne2$, $x\ne-1$. Porovnání exponentů dává $\frac{x+1}{x-2}=\frac{x-3}{x+1}$, tedy $(x+1)^2=(x-3)(x-2)$ a $7x=5$.
\item $x=\frac{5}{7}$ splňuje obě podmínky; $K=\{\frac{5}{7}\}$.
\end{enumerate}
\end{zadaniReseniBox}
\clearpage
\begin{zadaniReseniBox}[title={Domácí úkol}]
\begin{enumerate}[label=\alph*)]
\item Vyřeš $2^{3x-2}=2^{x+4}$.
\item Vyřeš $5^{x-1}=\frac{1}{25}$.
\item Vyřeš $3^{2x+5}=3^{x-1}$.
\item Vyřeš $\left(\frac{1}{5}\right)^{x+2}=\left(\frac{1}{5}\right)^{-1}$.
\item Vyřeš $7^{x^2-1}=7^{2x+2}$.
\item Vyřeš $2^{x^2}=2^{4x}$.
\item Rozhodni, zda $3^{x+2}=-9$ má reálný kořen.
\item Rozhodni, zda $2^{x+1}=0$ má reálný kořen.
\item Vysvětli, proč lze porovnat exponenty u $\left(\frac{3}{4}\right)^u=\left(\frac{3}{4}\right)^v$.
\item Najdi chybu v tvrzení: z $1^{x+3}=1^{2x-1}$ plyne $x=4$.
\end{enumerate}
\tcblower
\begin{enumerate}[label=\alph*)]
\item $3x-2=x+4$, tedy $x=3$.
\item $5^{x-1}=5^{-2}$, tedy $x=-1$.
\item $2x+5=x-1$, tedy $x=-6$.
\item $x+2=-1$, tedy $x=-3$.
\item $x^2-2x-3=0$, tedy $x=-1$ nebo $x=3$.
\item $x(x-4)=0$, tedy $x=0$ nebo $x=4$.
\item Nemá: levá strana je vždy kladná.
\item Nemá: mocnina se základem 2 je vždy kladná.
\item Funkce $\left(\frac{3}{4}\right)^x$ je přísně klesající, a tedy prostá.
\item Základ 1 není přípustný: rovnost $1=1$ platí pro všechna $x\in\mathbb R$.
\end{enumerate}
\end{zadaniReseniBox}
\poznamka{Domácí úkol asi 20--25 minut. Rozšíření 2 je dobrovolné; podmínky jmenovatelů zapisovat ještě před křížovým násobením.}
\end{document}
